\begin{papercard}{ferrari2026}
\cf{Stages} \stagetag{L} (LLM pipeline comparing two disclosure surfaces; model and prompt-decomposition comparison), \stagetag{D} (policy and Google Play Data Safety section captured per app), \stagetag{C} (Google's data-type taxonomy as shared vocabulary for both surfaces). arXiv preprint of April 2026 (Genoa and partners). Added from the gap search of 23 October 2026 as a recent measurement of how well an LLM compares a policy with a store disclosure, the comparison behind RQ4.

\cf{Problem and design} Since 2022 developers must fill in the Data Safety section, which should agree with the policy \pc{ferrari2026}{2}{In 2022, Google mandated the Data Safety Section (DSS) on Google Play, requiring developers to summarize their data practices}. PolicyGapper scrapes both, extracts data types from the policy with an LLM and reports types the policy declares but the Data Safety section omits, without app binaries \pc{ferrari2026}{2}{our stages: scraping, pre-processing, analysis, and post-processing, without requiring access to application binaries}. Google's terminology is used for both surfaces \pc{ferrari2026}{4}{Throughout this paper we adopt Google’s terminology for data types, data categories, and disclosure constraints}.

\cf{Models and prompts} Zero-shot prompting \pc{ferrari2026}{10}{leverage a Zero-shot Prompting strategy}; three models compared, one local open model, one cloud open model and one proprietary \pc{ferrari2026}{12}{We selected two freely accessible models and one paid model, namely Qwen3:8B, DeepSeek-V3.1:671B-Cloud, and Gemini 2.5 Pro}; four prompt decompositions from one unified prompt to per-component prompts were tested with Gemini 2.5 Pro. Gemini also read the policy as PDF, keeping tables and layout. Larger models could not be run locally \pc{ferrari2026}{12}{We attempted to run models with more than 8B parameters locally, but their computational cost did not scale}.

\cf{Results} 330 top-ranked apps from all 33 categories, Q3 2025; 2,689 omitted declarations (2,040 collection, 649 sharing) \pc{ferrari2026}{2}{We evaluate PolicyGapper on a dataset of 330 top-ranked apps spanning all 33 Google Play categories, collected in Q3 2025}. Validation on a stratified 10\,\% of policies, reviewed by hand, over three runs \pc{ferrari2026}{2}{Manual validation on a stratified 10\% subset, repeated across three independent runs, yields an average Precision of 0.75, Recall of 0.77, Accuracy of 0.69, and F1-score of 0.76}. Developers were notified afterwards.

\cf{Limits (card)} Preprint. Only omissions in one direction (policy declares, Data Safety omits) are tested; contradictions and vague wording are not. The validation subset is 33 policies, and the paper does not report agreement between human reviewers. Precision 0.75 means one in four reported gaps is wrong.

\cf{Use for this thesis} (1) Calibration for RQ4: an LLM comparison across two surfaces reached F1 0.76 with careful engineering, far below single-document classification (0.9+ in Xie et al.\ \pp{xie2025}). Cross-surface claims should be checked by hand, as the methodology already says for potential mismatches. (2) Mapping both surfaces onto one shared vocabulary (here Google's data types) is a precondition for comparison; the thesis needs the same for policy, listing, marketing and interface text. (3) Repeating the validation over three runs is a cheap stability check to copy. (4) Large-scale context: omissions between policy and Data Safety section are common, consistent with Arkalakis et al.\ \pp{arkalakis2024}.
\end{papercard}
